\documentclass[11pt]{article}
\usepackage{preamble}
% Auto-generated by figures/gen_figures.py. Do not edit by hand.

\newcommand{\Nresp}{374}
\newcommand{\Niv}{13}
\newcommand{\PctStudent}{81}
\newcommand{\PctDeveloper}{13}
\newcommand{\PctDailyWeekly}{82}
\newcommand{\NDaily}{147}
\newcommand{\PctUsedAI}{77}
\newcommand{\PctNeverAI}{10}
\newcommand{\PctReturnOften}{44}
\newcommand{\PctReturnNever}{7}
\newcommand{\PctMessyAgree}{59}
\newcommand{\MeanMessy}{3.6}
\newcommand{\PctTimeOverHour}{42}
\newcommand{\PctTimeOverThree}{14}
\newcommand{\PctNoTrustAuto}{36}
\newcommand{\PctDontKnowHow}{28}
\newcommand{\PctSetupCost}{24}
\newcommand{\NPrefA}{192}
\newcommand{\NPrefB}{182}
\newcommand{\PctPrefA}{51}
\newcommand{\PctWorkspaceTrust}{65}
\newcommand{\MeanWorkspaceTrust}{3.8}
\newcommand{\PctLocalFirst}{40}
\newcommand{\PctCloud}{22}
\newcommand{\PctNeedSpecifics}{38}
\newcommand{\PctKeepAfterMistake}{65}
\newcommand{\PctQuitAfterMistake}{10}
\newcommand{\PctLikelyUse}{47}
\newcommand{\PctUnlikelyUse}{22}
\newcommand{\MeanLikelyUse}{3.3}
\newcommand{\PctChurnHabit}{27}
\newcommand{\PctChurnMistake}{24}
\newcommand{\PctChurnRedundant}{20}
\newcommand{\PctChurnPrivacy}{17}
\newcommand{\PctChurnApproval}{10}
\newcommand{\FeatRankOneCount}{222}
\newcommand{\FeatRankTwoCount}{207}
\newcommand{\FeatRankThreeCount}{200}
\newcommand{\FeatRankFourCount}{191}
\newcommand{\FeatRankFiveCount}{154}
\newcommand{\FeatRankSixCount}{130}
\newcommand{\FeatRankSevenCount}{111}
\newcommand{\FeatRankEightCount}{105}
\newcommand{\FeatRankNineCount}{44}
\newcommand{\FeatTopName}{Resuming interrupted work}
\newcommand{\FeatTopCount}{222}
\newcommand{\FeatSpeechCount}{44}
\newcommand{\ThemeUndo}{17}
\newcommand{\ThemeLearn}{18}
\newcommand{\ThemeRedundant}{17}
\newcommand{\ThemeDataServer}{15}
\newcommand{\IvPositive}{10}
\newcommand{\IvKeepYes}{6}
\newcommand{\IvKeepConditional}{4}
\newcommand{\IvTrustShrug}{8}
\newcommand{\IvTrustStop}{2}
\newcommand{\PctPrefAHeavy}{44}
\newcommand{\PctPrefALight}{59}
\newcommand{\PctLocalStudent}{41}
\newcommand{\PctLocalDev}{36}


\title{\vspace{-1.5em}\textbf{Data Collection Findings and Synthesis}}
\author{}
\date{}

\begin{document}

% ======================================================================
% Title page
% ======================================================================
\begin{titlepage}
\centering
\vspace*{2.2cm}
{\Large ARTEMIS Foundational Research}\\[0.4cm]
{\LARGE\bfseries Data Collection Findings and Synthesis}\\[0.6cm]
{\large FAST NUCES Final-Year Project \textbf{F26-008}}\\[0.2cm]
{\small \textit{Adaptive Reasoning, Task Execution and Multi-agent Intelligence System}}\\[1.6cm]

\begin{tabular}{@{}ll@{}}
\toprule
Team member & Roll number \\
\midrule
Hashir Rauf & 23L-2572 \\
Hira Khalid & 23L-2594 \\
Doureesha Batool & 23L-2651 \\
\bottomrule
\end{tabular}

\vspace{0.9cm}

{\normalsize Supervisor: Mr.\ Mubasher Hussain}

\vfill
{\normalsize Internal research document. Not for external distribution.}\\[0.3cm]
{\normalsize Compiled \today}
\end{titlepage}

% ======================================================================
% Front matter
% ======================================================================
\pagenumbering{roman}
\setcounter{page}{1}
\tableofcontents
\vspace{1em}
\listoffigures
\clearpage
\pagenumbering{arabic}
\setcounter{page}{1}

% ======================================================================
\section{Executive Summary}
% ======================================================================

This report analyses the foundational research collected for ARTEMIS: a questionnaire with
\Nresp{} valid responses and \Niv{} coded in-person interviews. The purpose is to test whether
the problem ARTEMIS targets is real, whether the bounded-scope design is a genuine trust
advantage, and which features to build first.

The data confirms the problem. A large majority of respondents work across multiple files
often, lose meaningful time reorienting after a break, and see their file organisation decay
faster than they clean it. The reasons the work is not already automated are lack of trust in
automation and lack of know-how, not indifference. The bounded-workspace idea is supported:
roughly two thirds of respondents rate an assistant that only sees files they explicitly add
as meaningfully more trustworthy than one with full computer access, and local-first is the
single most common deployment preference. Predicted sustained use is moderate rather than
enthusiastic, and the most common churn risks are poor fit with existing habits and a mistake
the user cannot trust the tool to avoid again.

\subsubsection*{Headline findings}

The nine numbered findings (F1 to F9) are stated in context in the sections that follow. In
summary:
\begin{itemize}[leftmargin=3.2em,labelsep=0.6em]
  \item[\textbf{F1.}] The reorientation problem is real. \PctReturnOften{}\pctsym{} of
        respondents struggle to remember what they were doing when they return to a project
        after a break, and only \PctReturnNever{}\pctsym{} never do.
  \item[\textbf{F2.}] File disorder outpaces manual cleanup for \PctMessyAgree{}\pctsym{} of
        respondents (agree or strongly agree).
  \item[\textbf{F3.}] The problem costs real time at a moderate scale:
        \PctTimeOverHour{}\pctsym{} lose more than an hour a week.
  \item[\textbf{F4.}] The work is not automated because of distrust of automation
        (\PctNoTrustAuto{}\pctsym{}) and not knowing how (\PctDontKnowHow{}\pctsym{}), not
        because the tasks do not exist.
  \item[\textbf{F5.}] The bounded-workspace hypothesis holds. \PctWorkspaceTrust{}\pctsym{}
        rate a workspace-only assistant as more trustworthy than a full-access one
        (mean \MeanWorkspaceTrust{} on a 1 to 5 scale), and \PctLocalFirst{}\pctsym{} prefer
        local-first over cloud.
  \item[\textbf{F6.}] Users are forgiving of a single mistake. \PctKeepAfterMistake{}\pctsym{}
        would keep using ARTEMIS after one wrong suggestion; only \PctQuitAfterMistake{}\pctsym{}
        would abandon it outright.
  \item[\textbf{F7.}] The main retention risks are habit fit (\PctChurnHabit{}\pctsym{}) and an
        untrusted mistake (\PctChurnMistake{}\pctsym{}); approval fatigue is real but smaller
        (\PctChurnApproval{}\pctsym{}).
  \item[\textbf{F8.}] Feature demand is concentrated. Session resume, multi-document synthesis,
        and tidying a folder are the top three first-week features; speech input and output is
        last by a wide margin.
  \item[\textbf{F9.}] The interviews confirm the questionnaire and add effortless onboarding as
        a named condition for continued use.
\end{itemize}

% ======================================================================
\section{Research Method}
% ======================================================================

\subsection{Instruments}

The questionnaire had 17 items covering four areas: the respondent's context (role, how often
they do multi-file work, prior use of AI assistants); the problem (reorientation after a break,
file disorder, time lost, why repeated tasks are not automated); design and trust (approval-first
versus act-then-undo, the value of a bounded workspace, local versus cloud, tolerance for a
mistake, predicted sustained use, predicted reason to stop); and feature demand (a forced
top-three selection from ten candidate features, plus two free-text items). Likert items used a
1 to 5 scale.

The interview guide coded each conversation on five axes: the last time the participant lost
their place in multi-file work; the workaround they already use; their reaction to the ARTEMIS
concept; what a trust-breaking mistake would do to their usage; and whether they would keep
using the tool.

\subsection{Sample}

\begin{table}[H]
\centering
\caption{Response counts}
\begin{tabular}{@{}lr@{}}
\toprule
Instrument & Responses \\
\midrule
Questionnaire (valid rows) & \Nresp{} \\
In-person interviews (coded) & \Niv{} \\
\bottomrule
\end{tabular}
\end{table}

Recruitment was through FAST NUCES student and practitioner networks. The questionnaire ran
from 22 to 27 August 2026. The sample is student-heavy (\PctStudent{}\pctsym{} students,
\PctDeveloper{}\pctsym{} developers or software engineers, the remainder freelance or employed),
which matches the intended early-adopter profile for a final-year project but limits
generalisation.

\subsubsection{Limitations}
\begin{itemize}
  \item The sample is self-selected and drawn from a single institution's networks.
  \item Students dominate; the working-professional signal rests on a smaller subsample.
  \item ARTEMIS was described in text, not demonstrated. Predicted use and predicted churn are
        stated intentions, not observed behaviour.
  \item Free-text responses include a set of near-identical phrasings, which suggests some
        answer options were presented as suggested text. Theme counts should be read as
        indicative, not precise.
\end{itemize}

% ======================================================================
\section{Respondent Profile}
% ======================================================================

\begin{figure}[H]
\centering
\includegraphics[width=\textwidth]{fig_profile.pdf}
\caption{Respondent role and how often respondents do multi-file project work.}
\label{fig:profile}
\end{figure}

\subsection{Role and exposure to the problem domain}

As shown in Figure~\ref{fig:profile}, \PctStudent{}\pctsym{} of respondents are students and
\PctDeveloper{}\pctsym{} are developers or software engineers. Multi-file work is frequent:
\PctDailyWeekly{}\pctsym{} do it daily or a few times a week (\NDaily{} daily), and none of the
respondents do it less than a few times a month. The population that answered is therefore
squarely inside ARTEMIS's target use case.

\subsection{Prior experience with AI assistants}

\PctUsedAI{}\pctsym{} have used an AI assistant to help manage files, write code, or handle
tasks (regularly or occasionally), and only \PctNeverAI{}\pctsym{} have never tried one. The
audience is experienced enough to compare ARTEMIS against existing tools, which raises the bar
for differentiation (Section~\ref{sec:feedback}).

% ======================================================================
\section{Problem Validation}
% ======================================================================

\subsection{Reorientation after a break}

\begin{figure}[H]
\centering
\includegraphics[width=0.82\textwidth]{fig_return_struggle.pdf}
\caption{How often respondents return to a project after a break of a day or more and struggle
to remember what they were doing or what to do next.}
\label{fig:return}
\end{figure}

\PctReturnOften{}\pctsym{} of respondents answered \enquote{Often} or \enquote{Every time}
(Figure~\ref{fig:return}). Only \PctReturnNever{}\pctsym{} said it never happens. This is the
problem the first ARTEMIS feature, session resume, addresses directly.

\begin{finding}[The reorientation problem is real and widespread.]
Nearly half of respondents struggle to recover context after a break, and the overwhelming
majority experience it at least sometimes. Session resume targets a real, frequent pain point.
\end{finding}

\subsection{File disorder outpaces manual cleanup}

\begin{figure}[H]
\centering
\includegraphics[width=0.82\textwidth]{fig_messy_agree.pdf}
\caption{Agreement with \enquote{My folders and files tend to become messy or disorganized
faster than I can manually clean them up} (1 = strongly disagree, 5 = strongly agree).}
\label{fig:messy}
\end{figure}

\PctMessyAgree{}\pctsym{} of respondents rate this 4 or 5 (Figure~\ref{fig:messy}), mean
\MeanMessy{}. The tidy-a-folder feature has a clear constituency, though Section~\ref{sec:synth}
notes a countervailing segment that will not delegate organisation.

\begin{finding}[Most users cannot keep up with their own file organisation.]
A clear majority agree their files decay faster than they can clean them, which validates the
tidy-a-folder feature as a genuine need rather than a nice-to-have.
\end{finding}

\subsection{Self-reported time lost}

\begin{figure}[H]
\centering
\includegraphics[width=0.82\textwidth]{fig_time_lost.pdf}
\caption{Estimated total time lost per week to reorienting after breaks, tidying files,
repeating manual tasks, and searching scattered notes.}
\label{fig:time}
\end{figure}

\PctTimeOverHour{}\pctsym{} of respondents estimate losing more than an hour a week, and
\PctTimeOverThree{}\pctsym{} estimate more than three hours (Figure~\ref{fig:time}). The modal
answer is 30 to 60 minutes. This is a moderate, recurring cost rather than a crisis, which
matches the moderate predicted-use numbers in Section~\ref{sec:design}.

\begin{finding}[The problem costs real time, at a moderate scale.]
Most respondents lose 30 minutes to 3 hours a week. ARTEMIS should be positioned as removing
steady friction, not as solving an emergency.
\end{finding}

\subsection{Why the work is not already automated}

\begin{figure}[H]
\centering
\includegraphics[width=\textwidth]{fig_why_not_automated.pdf}
\caption{The honest reason respondents have not automated repeated manual file-handling tasks.}
\label{fig:whynot}
\end{figure}

The leading reasons are distrust of unsupervised automation (\PctNoTrustAuto{}\pctsym{}) and
not knowing how (\PctDontKnowHow{}\pctsym{}), followed by setup cost exceeding the payoff
(\PctSetupCost{}\pctsym{}). Only a minority say the tasks do not exist (Figure~\ref{fig:whynot}).

\begin{finding}[The barrier is trust and know-how, not absence of the need.]
Respondents want the work done but do not trust automation to do it unsupervised and do not
know how to set it up. An approval-gated assistant with zero configuration directly answers both
objections.
\end{finding}

% ======================================================================
\section{Design and Trust Signals}
\label{sec:design}
% ======================================================================

\subsection{Approval-first versus act-then-undo}

\begin{figure}[H]
\centering
\includegraphics[width=0.72\textwidth]{fig_assistant_pref.pdf}
\caption{Preference between Assistant A (asks for approval before every file change or outbound
action) and Assistant B (acts on reasonable requests automatically, with review and undo
afterward).}
\label{fig:pref}
\end{figure}

The split is close: \NPrefA{} for approval-first, \NPrefB{} for act-then-undo
(\PctPrefA{}\pctsym{} versus the remainder, Figure~\ref{fig:pref}).

\subsubsection{Implication}

There is no single right answer for the whole audience. Approval granularity should be a
user-configurable setting (for example: approve every action, approve by category, or trust the
session after the first few approvals), rather than a fixed product decision. This also
mitigates the approval-fatigue churn risk in Section~\ref{sec:churn}.

\subsection{The bounded-workspace hypothesis}

\begin{figure}[H]
\centering
\includegraphics[width=0.82\textwidth]{fig_workspace_trust.pdf}
\caption{Agreement with \enquote{An assistant that only sees files I explicitly add to a
workspace is meaningfully more trustworthy to me than one with full access to my computer}
(1 = strongly disagree, 5 = strongly agree).}
\label{fig:wstrust}
\end{figure}

\PctWorkspaceTrust{}\pctsym{} of respondents rate this 4 or 5, mean \MeanWorkspaceTrust{}
(Figure~\ref{fig:wstrust}). This is the most direct test in the survey of the project's central
hypothesis, that restricting an assistant's scope increases user trust, and the result supports
it.

\begin{finding}[The bounded workspace is a genuine trust advantage.]
Roughly two thirds of respondents find a workspace-only assistant meaningfully more trustworthy
than a full-access one. The bounded-scope design should be kept and made visible in the product,
not treated as an internal implementation detail.
\end{finding}

\subsection{Local-first versus cloud}

\begin{figure}[H]
\centering
\includegraphics[width=0.82\textwidth]{fig_local_vs_cloud.pdf}
\caption{Preference for a local, offline-first version versus a cloud-backed version.}
\label{fig:localcloud}
\end{figure}

Local-first is the single most common choice at \PctLocalFirst{}\pctsym{},
\PctNeedSpecifics{}\pctsym{} want more specifics before deciding, and only
\PctCloud{}\pctsym{} actively prefer cloud (Figure~\ref{fig:localcloud}). Privacy also appears
as a churn reason (\PctChurnPrivacy{}\pctsym{}, Section~\ref{sec:churn}) and as a recurring
free-text concern (Section~\ref{sec:feedback}). Respondents who prefer local-first also report
slightly higher predicted sustained use than those who prefer cloud, so the local-first option
is not costing engagement.

\subsection{Tolerance for a single mistake}

\begin{figure}[H]
\centering
\includegraphics[width=\textwidth]{fig_wrong_suggestion.pdf}
\caption{What respondents would do if ARTEMIS made one wrong or unwanted suggestion (for
example grouped files incorrectly or drafted an email with a mistake).}
\label{fig:mistake}
\end{figure}

\PctKeepAfterMistake{}\pctsym{} would keep using ARTEMIS after one mistake (either
\enquote{more cautious but keep using} or \enquote{one mistake would not change my usage}).
A further group would stop trusting only the specific feature that erred, and just
\PctQuitAfterMistake{}\pctsym{} would abandon the tool entirely (Figure~\ref{fig:mistake}).

\begin{finding}[One mistake is survivable if recovery is easy.]
Most users are forgiving of a single error, provided they can see and reverse it. This makes a
reliable, visible undo more important than eliminating every mistake.
\end{finding}

\subsection{Predicted sustained use}

\begin{figure}[H]
\centering
\includegraphics[width=0.82\textwidth]{fig_likely_use.pdf}
\caption{Likelihood of continuing to use ARTEMIS after the novelty wears off, if it worked
exactly as described (1 = very unlikely, 5 = very likely).}
\label{fig:likely}
\end{figure}

\PctLikelyUse{}\pctsym{} of respondents answered 4 or 5, \PctUnlikelyUse{}\pctsym{} answered
1 or 2, and the mean is \MeanLikelyUse{} (Figure~\ref{fig:likely}). This is a moderate,
plausible result for a described-but-not-demonstrated concept, and it sets a realistic
expectation: ARTEMIS should aim to convert the large \enquote{3} group through a good first
experience.

\subsection{Predicted reasons for churn}
\label{sec:churn}

\begin{figure}[H]
\centering
\includegraphics[width=\textwidth]{fig_stop_reason.pdf}
\caption{The single most likely reason a respondent would stop using ARTEMIS after a week.}
\label{fig:churn}
\end{figure}

The dominant risks are poor fit with existing habits (\PctChurnHabit{}\pctsym{}) and a mistake
the user cannot trust the tool to avoid again (\PctChurnMistake{}\pctsym{}), followed by
redundancy with what the user can already do (\PctChurnRedundant{}\pctsym{}) and privacy
(\PctChurnPrivacy{}\pctsym{}). Approval fatigue is real but the smallest named risk at
\PctChurnApproval{}\pctsym{} (Figure~\ref{fig:churn}).

\begin{finding}[Retention depends on habit fit and trustworthy recovery, not on fewer prompts.]
The two largest churn drivers are workflow fit and an untrusted mistake. Onboarding that plants
ARTEMIS inside an existing routine, plus visible undo, address the top risks. Reducing approval
prompts helps a smaller group and should not come at the cost of the trust the approval model
buys.
\end{finding}

% ======================================================================
\section{Feature Demand}
% ======================================================================

\subsection{First-week feature priorities}

\begin{figure}[H]
\centering
\includegraphics[width=\textwidth]{fig_features.pdf}
\caption{Number of times each feature was chosen in the forced top-three selection
(multi-select, \Nresp{} respondents).}
\label{fig:features}
\end{figure}

Ranked by selections (Figure~\ref{fig:features}):
\begin{enumerate}
  \item Resuming interrupted work (session resume), \FeatRankOneCount{}
  \item Pulling together notes and citations from multiple documents, \FeatRankTwoCount{}
  \item Tidying up a messy folder, \FeatRankThreeCount{}
  \item Code assistant (explain errors, suggest fixes as a diff), \FeatRankFourCount{}
  \item Detecting and offering to automate a repeated task, \FeatRankFiveCount{}
  \item Drafting documents from existing material, \FeatRankSixCount{}
  \item On-demand web search, \FeatRankSevenCount{}
  \item Drafting and sending emails with approval, \FeatRankEightCount{}
  \item Speech input and output, \FeatRankNineCount{}
\end{enumerate}

\begin{finding}[Demand is concentrated on the first three planned features; speech is last.]
Session resume, multi-document synthesis, and folder tidy are the clear top three. Speech input
and output is chosen roughly a fifth as often as the leader and should stay last in the build
order.
\end{finding}

\subsection{Alignment with the planned build order}

\begin{table}[H]
\centering
\caption{Planned ARTEMIS build order against measured first-week demand.}
\begin{tabular}{@{}>{\raggedright\arraybackslash}p{4.7cm}c>{\raggedright\arraybackslash}p{7.2cm}@{}}
\toprule
Planned feature (build order) & Demand rank & Comment \\
\midrule
Pick up where you left off (session resume) & 1 & Highest demand and lowest technical risk. Correct to build first. \\
Tidy up a folder & 3 & Strong demand; watch the \enquote{I organise my own files} segment. \\
Pull together notes & 2 & Second-highest demand. Bring forward if resources allow. \\
Notice a repeated task & 5 & Mid demand; pairs well with the automation-distrust finding (F4). \\
Code assistant & 4 & Solid demand, concentrated among the developer subsample. \\
Web search & 7 & Lower demand; keep on-demand only as planned. \\
Mailing & 8 & Lower demand; approval and logging are essential given trust concerns. \\
Document drafting & 6 & Mid demand. \\
Speech input and output & 9 & Lowest demand. Keep last. \\
\bottomrule
\end{tabular}
\end{table}

The measured demand order is close to the planned build order. The one adjustment worth
considering is prioritising multi-document synthesis, which ranks second in demand but is
currently scheduled after the file features.

\subsection{Unprompted requests}
\label{sec:feedback}

\begin{figure}[H]
\centering
\includegraphics[width=\textwidth]{fig_open_feedback_themes.pdf}
\caption{Recurring themes in free-text responses. Counts are indicative given repeated phrasings
in the raw data.}
\label{fig:themes}
\end{figure}

The recurring themes, by frequency (Figure~\ref{fig:themes}):
\begin{enumerate}
  \item \textbf{Learn my folder structure over time} (\ThemeLearn{}). The most requested
        capability not already on the roadmap. Users want ARTEMIS to remember how they organise
        work so they do not reconfigure it each session.
  \item \textbf{One-click universal undo} (\ThemeUndo{}). Named repeatedly as the thing that
        would make the tool trustworthy. Reinforces F5.
  \item \textbf{Already covered by Copilot or Claude Code} (\ThemeRedundant{}). A
        differentiation risk: an experienced audience is not sure what ARTEMIS adds. The answer
        is the bounded workspace and local-first operation, and this needs to be stated
        explicitly in all messaging.
  \item \textbf{Concern about how much data is sent to a server} (\ThemeDataServer{}).
        Supports the local-first default and a visible \enquote{What ARTEMIS Knows} panel.
  \item \textbf{File-consolidation requests} (auto-group files with similar names, merge
        duplicate versions, track the latest version across copies). Together these are the
        largest cluster and point to a \enquote{consolidate related files} capability within
        the tidy feature.
  \item \textbf{Work across cloud drives} (Google Drive, OneDrive), not only local folders.
        In tension with the local-first and bounded-scope design; note for later.
  \item \textbf{I organise my own files and will not delegate it}. A real segment that will not
        use the tidy feature regardless of quality.
\end{enumerate}

% ======================================================================
\section{Interview Findings}
% ======================================================================

\begin{figure}[H]
\centering
\includegraphics[width=\textwidth]{fig_interview.pdf}
\caption{Coded interview responses: reaction to the ARTEMIS concept, and whether the
participant would actually keep using it (\Niv{} interviews).}
\label{fig:interview}
\end{figure}

\subsection{The problem in participants' own words}

Almost every participant recognised the problem. The workarounds they already use are
consistent: ad hoc Markdown notes written by hand or with an AI agent to record where they left
off; manual, repeated restructuring because \enquote{every project has a different structure};
and simply tolerating the friction because automating it feels like more work than it saves.
Two participants said the problem does not affect them much, one because an IDE agent already
handles it.

\subsection{Reaction to the concept}

\IvPositive{} of \Niv{} participants reacted positively to the concept
(Figure~\ref{fig:interview}), with the rest split between conditional, mixed, and one negative.

\subsection{Trust-breaking moments}

Most participants (\IvTrustShrug{} of \Niv{}) said a single mistake would not stop them, though
several qualified this with \enquote{with caution}. A minority wanted the tool to learn from
being corrected. \IvTrustStop{} said a trust-breaking mistake would make them stop. This mirrors
the questionnaire result F5.

\subsection{Willingness to keep using}

\IvKeepYes{} said yes outright and \IvKeepConditional{} said they would keep using it
conditionally, with easy onboarding named as the most common condition, followed by privacy and
by the tool reaching sufficient maturity.

\begin{finding}[Interviews confirm the questionnaire and add onboarding as a named condition.]
Concept reception is positive, a single mistake is mostly survivable, and continued use is
conditional on onboarding being effortless.
\end{finding}

% ======================================================================
\section{Synthesis}
\label{sec:synth}
% ======================================================================

\subsection{What the data confirms}
\begin{itemize}
  \item The reorientation and file-disorder problems are real, frequent, and cost most users
        30 minutes to 3 hours a week (F1, F2, F3).
  \item The work is not automated because of distrust and lack of know-how, which an
        approval-gated, zero-configuration assistant is well placed to address (F4).
  \item The bounded workspace is a real trust advantage, and local-first is the most common
        deployment preference (F5).
  \item Users are forgiving of a single mistake when recovery is easy (F6), and the interviews
        agree (F9).
  \item Feature demand is concentrated on session resume, multi-document synthesis, and folder
        tidy; speech is last (F8).
\end{itemize}

\subsection{What the data challenges}
\begin{itemize}
  \item The audience is split almost evenly on approval-first versus act-then-undo, so a single
        fixed interaction model will not satisfy everyone.
  \item A visible segment states plainly that it organises its own files and will not delegate
        that, which caps the addressable audience for the tidy feature.
  \item An experienced audience is not sure ARTEMIS adds anything over Copilot or Claude Code.
        Differentiation is not self-evident and must be argued.
  \item Predicted sustained use is moderate (mean \MeanLikelyUse{}), not enthusiastic. The
        large neutral group is the one to win.
\end{itemize}

\subsection{Cross-cutting tensions}
\begin{itemize}
  \item \textbf{Approval friction versus trust.} The approval model is what makes the tool
        trustworthy, but approval fatigue is a named churn risk. Resolve with configurable
        granularity, not by weakening the model.
  \item \textbf{Local-first versus cross-drive convenience.} The privacy-driven local-first
        preference conflicts with requests to operate across Google Drive and OneDrive.
  \item \textbf{Personalisation versus bounded scope.} The top unmet request, learning the
        user's folder structure over time, implies persistent memory of the user's habits,
        which must be designed to stay inside the bounded-scope guarantee.
\end{itemize}

\subsubsection{Segment view}

% Auto-generated by figures/gen_figures.py. Do not edit by hand.

\begin{table}[H]
\centering
\caption{Selected measures by respondent segment.}
{\small
\begin{tabular}{@{}lrccc@{}}
\toprule
Segment & $n$ & Workspace-trust & Prefer & Prefer \\
 & & (mean, 1--5) & approval-first & local-first \\
\midrule
All respondents & 374 & 3.8 & 51\% & 40\% \\
Students & 302 & 3.8 & 51\% & 41\% \\
Developers / engineers & 47 & 3.7 & 51\% & 36\% \\
Regular AI-assistant users & 172 & 3.9 & 44\% & 40\% \\
Light or non AI-assistant users & 87 & 3.8 & 59\% & 41\% \\
\bottomrule
\end{tabular}}
\end{table}


Students and developers give near-identical workspace-trust ratings and identical approval-first
shares, so role does not predict the interaction preference. Experience does move it: regular
AI-assistant users prefer approval-first less often than light or non-users
(\PctPrefAHeavy{}\pctsym{} versus \PctPrefALight{}\pctsym{}), which is consistent with heavier
users being more comfortable delegating. The effect is modest and neither group is close to
unanimous, so the practical conclusion is unchanged: the interaction model should be a setting
rather than a fixed decision. Local-first preference is broadly consistent across all segments
(\PctLocalDev{}\pctsym{} to \PctLocalStudent{}\pctsym{}).

% ======================================================================
\section{Recommendations for the Build}
% ======================================================================

\begin{enumerate}
  \item \textbf{Build session resume, multi-document synthesis, and folder tidy first; keep
        speech input and output last.} Motivated by F8. Consider moving multi-document synthesis
        ahead of the later file features given it ranks second in demand.
  \item \textbf{Ship a reliable, visible, one-click undo before any outward-action feature.}
        Motivated by F6 and the free-text requests. Undo is the single most cited condition for
        trusting the tool.
  \item \textbf{Make onboarding effortless and design it to plant ARTEMIS inside an existing
        routine.} Motivated by F7 and F9, which name onboarding as the top condition for
        continued use.
  \item \textbf{Keep the bounded opt-in workspace and offer a local-first mode, and make both
        visible in the product.} Motivated by F5 and the privacy churn and free-text signals.
        This is also the differentiation answer to the \enquote{already covered by Copilot}
        objection.
  \item \textbf{Make approval granularity user-configurable} (per action, per category, or
        session trust after the first approvals). Motivated by the near-even A versus B split
        and the approval-fatigue churn risk.
  \item \textbf{Add \enquote{learn my workspace structure over time} to the roadmap.} It is the
        top requested capability not currently planned. Design it as per-workspace memory of the
        user's organisation preferences so it stays within the bounded-scope guarantee.
  \item \textbf{State the differentiation from Copilot and Claude Code explicitly} in the
        proposal, the UI, and any external material. The target audience already uses those
        tools and will not infer the difference.
  \item \textbf{Treat file consolidation as a first-class part of the tidy feature}
        (group files with similar names, flag duplicate versions, identify the latest version),
        since together these are the largest free-text cluster.
\end{enumerate}

\subsection{Open questions carried forward}
\begin{itemize}
  \item How to expose configurable approval granularity without making the first-run experience
        complex.
  \item Whether and how to support cloud drives without breaking the local-first and
        bounded-scope guarantees.
  \item How persistent \enquote{learn my structure} memory is represented in the
        \enquote{What ARTEMIS Knows} panel.
  \item How to convert the large neutral group (predicted-use rating 3) through the first
        session. These extend the existing internal open-questions note.
\end{itemize}

% ======================================================================
\appendix
\clearpage
\section{Appendix}
% ======================================================================

\subsection{Questionnaire items}
\begin{enumerate}
  \item Timestamp.
  \item Email address.
  \item How often do you work on projects that involve managing multiple files, documents, or
        folders (coursework, code projects, client deliverables, research)?
  \item What best describes you?
  \item How often do you return to a project after a break of a day or more and struggle to
        remember what you were doing or what to do next?
  \item \enquote{My folders and files tend to become messy or disorganized faster than I can
        manually clean them up} (1 = strongly disagree, 5 = strongly agree).
  \item If you repeat the same manual file-handling task often, why have you not automated it?
        Pick the honest reason.
  \item Roughly how much total time per week do you think you lose to these issues (reorienting
        after breaks, tidying files, repeating manual tasks, searching scattered notes)?
  \item Have you used an AI assistant (Copilot, ChatGPT, Gemini, a computer-use agent) to help
        manage files, write code, or handle tasks on your computer?
  \item Which of these two assistants would you actually prefer to use day to day? (A: asks for
        approval before every file change or outbound action; B: acts on reasonable requests
        automatically, with review and undo afterward.)
  \item \enquote{An assistant that only sees files I explicitly add to a workspace is
        meaningfully more trustworthy to me than one with full access to my computer}
        (1 = strongly disagree, 5 = strongly agree).
  \item Suppose ARTEMIS made a wrong or unwanted suggestion once (grouped your files
        incorrectly, or drafted an email with a mistake). Would you keep using it?
  \item Which of the following features would you actually use in your first week, if it worked
        exactly as described? Select your top 3.
  \item A local, offline-first version versus a cloud-backed version. Which would you actually
        choose?
  \item If ARTEMIS worked exactly as described, how likely would you be to actually use it for
        your own work, not \enquote{try once} but keep using after the novelty wears off?
        (1 = very unlikely, 5 = very likely).
  \item What is the single most likely reason you would stop using ARTEMIS after a week, if you
        had to guess right now?
  \item Is there anything about how you currently manage multi-file work, a frustration, a
        workaround, or a feature you wish existed, that this questionnaire did not ask about but
        that we should know before building a tool like this?
\end{enumerate}

\subsection{Interview coding scheme}
\begin{table}[H]
\centering
\caption{Interview coding axes.}
\begin{tabular}{@{}>{\raggedright\arraybackslash}p{3.6cm}>{\raggedright\arraybackslash}p{10.3cm}@{}}
\toprule
Axis & Coded values \\
\midrule
The last time you lost your place & Problem exists / happened a lot / never happened / other free-text descriptions. \\
The workaround you already have & Markdown notes, manual restructuring, \enquote{too lazy to automate}, not equipped, no need. \\
Reacting to the concept & Positive / semi-positive / mixed / negative. \\
The trust-breaking moment & Shrug it off (sometimes with caution) / conditional / needs human feedback / stop using. \\
Would you actually keep using it & Yes / conditional (onboarding, maturity) / privacy concern / no. \\
\bottomrule
\end{tabular}
\end{table}

\subsection{Complete response distributions}

Counts and percentages for every closed questionnaire item and every coded interview axis.

% Auto-generated by figures/gen_figures.py. Do not edit by hand.

\begin{table}[H]
\centering
\caption{Q4. What best describes you?}
\begin{tabular}{@{}lrr@{}}
\toprule
Response & Count & Percent \\
\midrule
Student & 302 & 80.7\% \\
Developer / Software Engineer & 47 & 12.6\% \\
Freelance & 12 & 3.2\% \\
Employed & 7 & 1.9\% \\
Vibe coder & 6 & 1.6\% \\
\midrule
Total & 374 & 100.0\% \\
\bottomrule
\end{tabular}
\end{table}

\begin{table}[H]
\centering
\caption{Q3. Frequency of multi-file project work}
\begin{tabular}{@{}lrr@{}}
\toprule
Response & Count & Percent \\
\midrule
Daily & 147 & 39.3\% \\
A Few Times a Week & 159 & 42.5\% \\
A Few Times a Month & 68 & 18.2\% \\
\midrule
Total & 374 & 100.0\% \\
\bottomrule
\end{tabular}
\end{table}

\begin{table}[H]
\centering
\caption{Q9. Prior use of an AI assistant for files / code / tasks}
\begin{tabular}{@{}lrr@{}}
\toprule
Response & Count & Percent \\
\midrule
Yes, regularly & 172 & 46.0\% \\
Yes, occasionally & 115 & 30.7\% \\
Tried it once or twice & 48 & 12.8\% \\
Never & 39 & 10.4\% \\
\midrule
Total & 374 & 100.0\% \\
\bottomrule
\end{tabular}
\end{table}

\begin{table}[H]
\centering
\caption{Q5. Return after a break and struggle to remember what to do next}
\begin{tabular}{@{}lrr@{}}
\toprule
Response & Count & Percent \\
\midrule
Never & 26 & 7.0\% \\
Rarely & 69 & 18.4\% \\
Sometimes & 113 & 30.2\% \\
Often & 137 & 36.6\% \\
Every time & 29 & 7.8\% \\
\midrule
Total & 374 & 100.0\% \\
\bottomrule
\end{tabular}
\end{table}

\begin{table}[H]
\centering
\caption{Q6. Files become messy faster than I can clean them up (1 to 5)}
\begin{tabular}{@{}lrr@{}}
\toprule
Response & Count & Percent \\
\midrule
1 & 9 & 2.4\% \\
2 & 44 & 11.8\% \\
3 & 100 & 26.7\% \\
4 & 145 & 38.8\% \\
5 & 76 & 20.3\% \\
\midrule
Total & 374 & 100.0\% \\
\bottomrule
\end{tabular}
\end{table}

\begin{table}[H]
\centering
\caption{Q8. Estimated time lost per week to these issues}
\begin{tabular}{@{}lrr@{}}
\toprule
Response & Count & Percent \\
\midrule
Under 30 minutes & 94 & 25.1\% \\
30 minutes - 1hour & 122 & 32.6\% \\
1-3 hours & 106 & 28.3\% \\
3-5 hours & 26 & 7.0\% \\
Over 5 Hours & 26 & 7.0\% \\
\midrule
Total & 374 & 100.0\% \\
\bottomrule
\end{tabular}
\end{table}

\begin{table}[H]
\centering
\caption{Q7. Why repeated manual file tasks are not automated}
\begin{tabular}{@{}lrr@{}}
\toprule
Response & Count & Percent \\
\midrule
I don't trust automation to do it correctly without supervision & 133 & 35.8\% \\
I don't know how / it feels out of reach & 105 & 28.2\% \\
It would take longer to set up than the time it saves & 88 & 23.7\% \\
Not applicable ; I don't have repeated manual tasks & 46 & 12.4\% \\
\midrule
Total & 372 & 100.0\% \\
\bottomrule
\end{tabular}
\end{table}

\begin{table}[H]
\centering
\caption{Q10. Preferred assistant: approve-first (A) vs act-then-undo (B)}
\begin{tabular}{@{}lrr@{}}
\toprule
Response & Count & Percent \\
\midrule
Assistant A
(approve first) & 192 & 51.3\% \\
Assistant B
(act, then undo) & 182 & 48.7\% \\
\midrule
Total & 374 & 100.0\% \\
\bottomrule
\end{tabular}
\end{table}

\begin{table}[H]
\centering
\caption{Q11. A workspace-only assistant is more trustworthy than a full-access one (1 to 5)}
\begin{tabular}{@{}lrr@{}}
\toprule
Response & Count & Percent \\
\midrule
1 & 14 & 3.7\% \\
2 & 34 & 9.1\% \\
3 & 82 & 21.9\% \\
4 & 121 & 32.4\% \\
5 & 123 & 32.9\% \\
\midrule
Total & 374 & 100.0\% \\
\bottomrule
\end{tabular}
\end{table}

\begin{table}[H]
\centering
\caption{Q14. Local-first vs cloud-backed preference}
\begin{tabular}{@{}lrr@{}}
\toprule
Response & Count & Percent \\
\midrule
Local First even if its slower or weaker & 150 & 40.1\% \\
I would need to know more specifics before deciding either way & 143 & 38.2\% \\
Cloud-backed even if its less private & 81 & 21.7\% \\
\midrule
Total & 374 & 100.0\% \\
\bottomrule
\end{tabular}
\end{table}

\begin{table}[H]
\centering
\caption{Q12. Reaction to one wrong or unwanted suggestion}
\begin{tabular}{@{}lrr@{}}
\toprule
Response & Count & Percent \\
\midrule
I'd be more cautious but keep using it & 186 & 49.7\% \\
I'd stop trusting its suggestions in that specific feature  area & 92 & 24.6\% \\
Yes, one mistake wouldn't change my usage & 57 & 15.2\% \\
I'd likely stop using the tool altogether & 39 & 10.4\% \\
\midrule
Total & 374 & 100.0\% \\
\bottomrule
\end{tabular}
\end{table}

\begin{table}[H]
\centering
\caption{Q15. Likelihood of sustained use if it worked as described (1 to 5)}
\begin{tabular}{@{}lrr@{}}
\toprule
Response & Count & Percent \\
\midrule
1 & 27 & 7.2\% \\
2 & 56 & 15.0\% \\
3 & 115 & 30.7\% \\
4 & 114 & 30.5\% \\
5 & 62 & 16.6\% \\
\midrule
Total & 374 & 100.0\% \\
\bottomrule
\end{tabular}
\end{table}

\begin{table}[H]
\centering
\caption{Q16. Single most likely reason to stop using ARTEMIS after a week}
\begin{tabular}{@{}lrr@{}}
\toprule
Response & Count & Percent \\
\midrule
I forget to use it/ it doesn't fit into my existing habit & 102 & 27.3\% \\
It makes a mistake I dont trust it to avoid again & 89 & 23.8\% \\
It doesn't do anything I couldn't already do myself easily & 74 & 19.8\% \\
Privacy or data-control concerns & 62 & 16.6\% \\
It asks for approval too often and slows me down & 37 & 9.9\% \\
I don't think I would stop using it & 10 & 2.7\% \\
\midrule
Total & 374 & 100.0\% \\
\bottomrule
\end{tabular}
\end{table}

\begin{table}[H]
\centering
\caption{Q13. Features that would actually be opened in the first week (pick top 3)}
\begin{tabular}{@{}lrr@{}}
\toprule
Response & Count & Percent \\
\midrule
Resuming interrupted work (session resume) & 222 & 16.3\% \\
Pulling together notes/citations from multiple documents & 207 & 15.2\% \\
Tidying up a messy folder & 200 & 14.7\% \\
Code assistant (explain errors, suggest fixes as a diff) & 191 & 14.0\% \\
Detecting and offering to automate a repeated task & 154 & 11.3\% \\
Drafting documents from existing material & 130 & 9.5\% \\
On-demand web search & 111 & 8.1\% \\
Drafting and sending emails (with approval) & 105 & 7.7\% \\
Speech input/output & 44 & 3.2\% \\
\midrule
Total & 1364 & 100.0\% \\
\bottomrule
\end{tabular}
\end{table}

\begin{table}[H]
\centering
\caption{Q17. Recurring themes in open-ended feedback}
\begin{tabular}{@{}lrr@{}}
\toprule
Response & Count & Percent \\
\midrule
Learn my folder structure over time & 18 & 14.8\% \\
One-click universal undo & 17 & 13.9\% \\
Already covered by Copilot / Claude Code & 17 & 13.9\% \\
Concern over data sent to a server & 15 & 12.3\% \\
Auto-group files with similar names & 14 & 11.5\% \\
Merge duplicate versions of a file & 11 & 9.0\% \\
Track the latest version across copies & 10 & 8.2\% \\
Work across cloud drives (Drive, OneDrive) & 10 & 8.2\% \\
I organise my own files and will not delegate it & 10 & 8.2\% \\
\midrule
Total & 122 & 100.0\% \\
\bottomrule
\end{tabular}
\end{table}

\begin{table}[H]
\centering
\caption{Interview: reaction to the concept}
\begin{tabular}{@{}lrr@{}}
\toprule
Response & Count & Percent \\
\midrule
Positive & 10 & 76.9\% \\
Semi-positive & 1 & 7.7\% \\
Mixed & 1 & 7.7\% \\
Negative & 1 & 7.7\% \\
\midrule
Total & 13 & 100.0\% \\
\bottomrule
\end{tabular}
\end{table}

\begin{table}[H]
\centering
\caption{Interview: would you actually keep using it}
\begin{tabular}{@{}lrr@{}}
\toprule
Response & Count & Percent \\
\midrule
Yes & 6 & 46.2\% \\
Conditional & 4 & 30.8\% \\
Privacy concern & 1 & 7.7\% \\
No & 2 & 15.4\% \\
\midrule
Total & 13 & 100.0\% \\
\bottomrule
\end{tabular}
\end{table}

\begin{table}[H]
\centering
\caption{Interview: reaction to a trust-breaking moment}
\begin{tabular}{@{}lrr@{}}
\toprule
Response & Count & Percent \\
\midrule
Shrug it off & 8 & 61.5\% \\
Conditional & 2 & 15.4\% \\
Stop using & 2 & 15.4\% \\
Needs human feedback & 1 & 7.7\% \\
\midrule
Total & 13 & 100.0\% \\
\bottomrule
\end{tabular}
\end{table}



\end{document}
